\subsection{单位根}

定义2. --- 元素 \(\zeta\) （属于 \(\mathbf{K}\) ）称为单位根，如果存在整数 \(n > 0\) 使得 \(\zeta^n = 1\) ；对每个整数 \(n > 0\) （使得 \(\zeta^n = 1\) ）， \(\zeta\) 被称为 \(n\) 次单位根。

这等同于说，单位根是 K 的非零元素所成乘法群 \(K^{*}\) 中有限阶的元素（I，第51页）。单位根构成一个子群 \(\mu_{\infty}(\mathbf{K})\) （ \(K^{*}\) 的），n 次单位根构成一个子群 \(\mu_{n}(\mathbf{K})\) （ \(\mu_{\infty}(\mathbf{K})\) 的）。我们有 \(\mu_{\infty}(\mathbf{K}) = \bigcup_{n \geqslant 1} \mu_{n}(\mathbf{K})\) ，且 \(\mu_{n}(\mathbf{K}) \subset \mu_{m}(\mathbf{K})\) ，如果 n 整除

\begin{enumerate}
\def\labelenumi{\alph{enumi}.}
\setcounter{enumi}{12}
\tightlist
\item
  对每一个单位根 \(\zeta\) ，存在一个最小的整数 \(n \geqslant 1\) 使得 \(\zeta\) 属于 \(\mu_{n}(\mathbf{K})\) ，即 \(\zeta\) 在群 \(K^{*}\) 中的阶。
\end{enumerate}

群 \(\mu_n(\mathbf{K})\) 是多项式 \(\mathbf{X}^n - 1\) 的根所成的集合，因此是阶 \(\leqslant n\) 的有限群。设 \(p\) 为 \(\mathbf{K}\) 的特征。当 \(p = 0\) ，或者当 \(p \neq 0\) 且 \(n\) 不被 \(p\) 整除时，导数 \(n\mathbf{X}^{n-1}\) （ \(\mathbf{X}^n - 1\) 的导数）与 \(\mathbf{X}^n - 1\) 互素，因此多项式 \(\mathbf{X}^n - 1\) 是可分的；此外若 \(\mathbf{K}\) 是代数闭的，则 \(\mu_n(\mathbf{K})\) 是一个由 \(n\) 个元素组成的群。

假设 K 的特征为非零的 p，并令 \(r \geqslant 0\) 为一个整数；由于映射 \(x \mapsto x^{p^{r}}\) 从 K 到 K 是单射，我们有 \(\mu_{np^{r}}(\mathbf{K}) = \mu_n(\mathbf{K})\) 对于每个整数 \(n \geqslant 1\) 。

我们注意到，一个域可能不含除1以外的任何 \(n\) 次单位根：例如素域 \(\mathbf{Q}\) 和 \(\mathbf{F}_2\) 对任何奇整数 \(n\) 就是这种情形。

定理1. --- 设 p 为 K 的特征指数，且设 n \textgreater{} 0 是一个整数。则 K 中 n 次单位根所成的群 \(\mu_{n}(\mathbf{K})\) 是循环的，且其阶整除 n；当 K 是代数闭的且 n 与 p 互素时，群 \(\mu_{n}(\mathbf{K})\) 是 n 阶循环群。

只需证明定理的第一个断言，它是以下更精确的引理的推论：

引理1. --- 设 G 是 \(\mathbf{K}^*\) 的一个阶为 \(m\) 的有限子群；则 G 是循环群，并且我们有 \(\mathrm{G} = \mu_m(\mathrm{K})\) 。

将 G 视为 Z-模；对所有 \(x \in G\) 有 mx = 0，因此 G 的零化子是形如 rZ 的理想，其中整数 \(r \geqslant 1\) 整除 m。因此我们有 \(G \subset \mu_{r}(K)\) 。由引理 4（V，第 74 页）应用于 A = Z 和 M = G，存在 G 中一个阶为 r 的元素 x；设 \(G'\) 为 G 中由 x 生成的循环子群。我们有 \(G' \subset \mu_{r}(K)\) ，Card \((G') = r\) 与 Card \(\mu_{r}(K) \leqslant r\) ；因此我们有

\(G' = \mu_r(K) \supset G\) ，且 \(G\) 是阶为 \(r\) 的循环群，从而等于 \(\mu_r(K)\) 。由于 \(G\) 的阶为 \(m\) ，我们有 \(m = r\) ，引理得证。

命题2. --- 设 K 是代数闭的，并设 p 为其特征指数。存在 \(\mu_{\infty}(\mathbf{K})\) 到群 \(\mathbf{Q}/\mathbf{Z}[1/p]\) 上的一个同构。

我们用 \(\mathbf{Z}[1/p]\) 表示 \(\mathbf{Q}\) 的由 \(1/p\) 生成的子环，即形如 \(a/p^n\) （ \(a \in \mathbf{Z}\) ， \(n \geqslant 1\) ）的有理数组成的集合；因此我们有 \(\mathbf{Z}[1/p] = \mathbf{Z}\) （当 \(\mathbf{K}\) 的特征为0时）。

设 \((\nu_{n})_{n\geqslant 1}\) 为由所有不被 \(p\) 整除的整数组成的严格递增序列（若 \(p\neq 1\) ）；令 \(\lambda_{n} = \nu_{1}\nu_{2}\dots \nu_{n}\) ，并用 \(\mathrm{H}_n\) 表示 \(\lambda_{n}\) 次单位根所成的群；我们有 \(\mathrm{H}_{n + 1}\supset \mathrm{H}_n\) ，且 \(\mu_{\infty}(\mathbf{K}) = \cup_{n}\mathrm{H}_{n}\) 。由于 \(\mathrm{H}_n\) 是阶为 \(\lambda_{n}\) 的循环群（定理1），存在一个序列

\((\alpha_{n})_{n\geqslant 1}\) ，由单位根组成，使得 \(\alpha_{n}\) 生成 \(\mathrm{H}_n\) ，且 \(\alpha_{n} = \alpha_{n + 1}^{\nu_{n + 1}}\)

进一步设 \(\beta_{n}\) 为模 \(\mathbf{Z}[1/p]\) 的 \(1/\lambda_{n}\) 的剩余类，并令 \(\mathbf{H}_{n}^{\prime}\) 为 \(\mathbf{Q}/\mathbf{Z}[1/p]\) 的由 \(\beta_{n}\) 生成的循环子群。显然有 \(\beta_{n} = \nu_{n+1}\beta_{n+1}\) ，且 \(\mathbf{H}_{n}^{\prime}\) 的阶为 \(\lambda_{n}\) ，因为 \(\lambda_{n}\) 不被 \(p\) 整除（当 \(p \neq 1\) 时）。于是对所有 \(n \geqslant 1\) 存在同构 \(\varphi_{n}: \mathrm{H}_{n} \to \mathrm{H}_{n}^{\prime}\) 使得 \(\varphi_{n}(\alpha_{n}) = \beta_{n}\) ，且关系式 \(\alpha_{n} = \alpha_{n+1}^{\nu_{n+1}}\) ， \(\beta_{n} = \nu_{n+1}\beta_{n+1}\) 表明 \(\varphi_{n+1}\) 延拓 \(\varphi_{n}\) 。最后，存在唯一的同构 \(\varphi\) （从 \(\mu_{\infty}(\mathbf{K})\) 到 \(\mathbf{Q}/\mathbf{Z}[1/p]\) 上），它延拓诸同构 \(\varphi_{n}\) ，即 \(\varphi(\alpha_{n}) = \beta_{n}\) 对所有 \(n \geqslant 1\) 成立。

注记. --- 1) 当 K 是特征为0的代数闭域时，群 \(\mu_{\infty}(\mathbf{K})\) 因此（非典范地）同构于 \(\mathbf{Q}/\mathbf{Z}\) 。* 当 K 是复数域 \(\mathbf{C}\) 时，我们可以写出这样一个明确的同构；实际上，映射 \(x \mapsto e^{2\pi ix}\) 是 \(\mathbf{Q}\) 到 \(\mathbf{C}^*\) 中的一个同态，其核为 \(\mathbf{Z}\) ，其像为 \(\mu_{\infty}(\mathbf{C})\) ；因此通过过渡到商群，它定义了 \(\mathbf{Q}/\mathbf{Z}\) 到 \(\mu_{\infty}(\mathbf{C})\) 上的一个同构。*

\begin{enumerate}
\def\labelenumi{\arabic{enumi})}
\setcounter{enumi}{1}
\item
  下述结果可以证明（参见 V，第165页，习题 21）：设 G 和 H 为两个交换群，其所有元素均为有限阶。假设对每个整数 \(n \geqslant 1\) ，方程 \(nx = 0\) 在 G 和 H 中具有相同个数的解（设为有限）。则群 G 与 H 同构。这给出了命题2的一个新证明。
\item
  对每个素数 \(l\) ，令 \(\mu_{l^{\infty}}(\mathbf{K}) = \bigcup_{n\geqslant 0}\mu_{l^{n}}(\mathbf{K})\) 。当 \(l\) 是
\end{enumerate}

\(p\) （ \(\mathbf{K}\) 的特征）时，我们有 \(\mu_{p^{\infty}}(\mathbf{K}) = \{1\}\) 。由 I，第80页，定理4，我们推出 \(\mu_{\infty}(\mathbf{K})\) 是诸子群 \(\mu_{l^{\infty}}(\mathbf{K})\) （其中 \(l\) 遍历所有异于 \(p\) 的素数）的直和。对给定的素数 \(l\) ，只有两种可能的情形：或者 \(\mu_{l^{\infty}}(\mathbf{K})\) 是有限的，此时 \(\mu_{l^{\infty}}(\mathbf{K})\) 同构于 \(\mathbf{Z}/l^{n}\mathbf{Z}\) ，对某个 \(n\) （定理1）；或者 \(\mu_{l^{\infty}}(\mathbf{K})\) 是无限的，此时 \(\mu_{l^{\infty}}(\mathbf{K})\) 同构于 \(\mathbf{Z}[l^{-1}]/\mathbf{Z}\) （参见注记2）。

